\documentclass[12pt]{article}
\usepackage{amsmath, amssymb}
\usepackage{geometry}
\usepackage{enumitem}
\usepackage{tcolorbox}
\usepackage{titlesec}

\geometry{margin=1in}

\tcbuselibrary{skins, breakable}
\newtcolorbox{answerbox}{
    colback=gray!8,
    colframe=gray!50,
    boxrule=0.5pt,
    arc=4pt,
    title=Answer,
    fonttitle=\bfseries
}

\title{\textbf{ESE 650: Learning in Robotics}\\[0.5em]
Problem 4 --- Variant Questions}
\date{}

\begin{document}
\maketitle

\noindent
The following three questions are variations of Problem 4 from the midterm,
covering LiDAR point cloud accumulation and NeRF-based 3D reconstruction.

\bigskip
\hrule
\bigskip

%% -------------------------------------------------------
%% Variant 1
%% -------------------------------------------------------
\section*{Variant 1 (Basic)}

The original point cloud accumulation formula is
\[
    P = \left\{ x_i + d_i^j \,\omega_i^j \;:\; i = 1,\ldots,m,\; j = 1,\ldots,n \right\}.
\]
Suppose the LiDAR sensor has a systematic error: every distance measurement is
scaled by a known factor $s > 0$, i.e., the true distance is $s \cdot d_i^j$.
Write the corrected point cloud accumulation formula.

\begin{answerbox}
Because the error affects the distance measurement $d_i^j$ and not the sensor
position $x_i$, we simply replace $d_i^j$ with $s \cdot d_i^j$:
\[
    P = \left\{ x_i + s\, d_i^j \,\omega_i^j \;:\; i = 1,\ldots,m,\; j = 1,\ldots,n \right\}.
\]
Note that $s$ is a \emph{multiplicative} scaling factor, not an additive offset.
Writing $d_i^j + s$ would be physically incorrect.
\end{answerbox}

\bigskip
\hrule
\bigskip

%% -------------------------------------------------------
%% Variant 2
%% -------------------------------------------------------
\section*{Variant 2 (Intermediate)}

The original method uses a NeRF to complete holes in the accumulated point cloud,
assuming a \emph{static} scene. Now suppose the scene contains \textbf{dynamic
objects} (e.g., moving vehicles), and the $m$ LiDAR scans are captured at
different times.

\begin{enumerate}[label=(\alph*)]
    \item What problem arises when applying the original NeRF method directly?
    \item Propose two approaches to address this issue---one based on
          \emph{data preprocessing} and one based on \emph{model modification}.
\end{enumerate}

\begin{answerbox}
\textbf{(a) Problem.}
NeRF assumes a static scene and fits a single neural radiance field to all
scans jointly. For a location that is occupied by a vehicle at some times and
empty at others, the network averages the occupancy across time steps,
learning an incorrect intermediate density---the reconstructed surface is
neither a solid vehicle nor empty space, resulting in \emph{ghosting} or
blurring artifacts.

\medskip
\textbf{(b) Solutions.}

\textit{Data preprocessing.}
Before training, filter out \emph{inconsistent points}: any spatial location
that is occupied in some scans but empty in others is flagged as dynamic and
removed. Only points that are consistently present across all time steps are
kept and fed to NeRF.

\textit{Model modification.}
Extend NeRF with a time dimension $t$, i.e., the network takes input
$(x, y, z, \theta, \phi, t)$. A common architecture (e.g., D-NeRF) adds a
\emph{deformation field} $\Delta x = f(x,y,z,t)$ that maps each query point
at time $t$ back to a canonical static space, where a standard NeRF is
evaluated. This allows the model to represent moving objects explicitly.
\end{answerbox}

\bigskip
\hrule
\bigskip

%% -------------------------------------------------------
%% Variant 3
%% -------------------------------------------------------
\section*{Variant 3 (Open-ended Design)}

Relax the assumption that ``viewpoints are rich enough.''  Suppose all $m$
LiDAR scans are taken from \textbf{the same side} (e.g., only the left side of
a building), so the right side has \emph{zero} point cloud coverage.

\begin{enumerate}[label=(\alph*)]
    \item Can the original NeRF-based method reliably complete the right side?
          Explain why or why not.
    \item Propose approaches to address this limitation without requiring
          additional LiDAR scans from the right side (and one approach that
          does use additional sensors).
\end{enumerate}

\begin{answerbox}
\textbf{(a) NeRF cannot reliably complete the unseen side.}
NeRF's hole-filling ability comes from the neural network's \emph{interpolation}
capacity: given dense observations surrounding a small gap, it can infer the
missing region. However, the right side requires \emph{extrapolation}---there
are no observations at all on that side. Without any training signal,
the network's density predictions on the right side are essentially arbitrary
(typically near zero), and cannot be trusted.

\medskip
\textbf{(b) Solutions.}

\textit{Symmetry prior (no extra sensors).}
If the structure is approximately symmetric, mirror the left-side point cloud
to generate pseudo-observations for the right side. Adding small random
perturbations (noise) to the mirrored points prevents over-reliance on perfect
symmetry. This is a form of data augmentation, but fails when the two sides
differ structurally (e.g., an entrance only on the right).

\textit{Shape completion model (no extra sensors).}
Use a deep learning model pre-trained on large 3D shape datasets (shape
completion networks) to infer the missing geometry from partial observations.
These models encode learned structural priors about how objects typically look,
enabling \emph{informed} extrapolation rather than blind guessing.

\textit{Additional sensor (extra LiDAR or camera).}
Place a second LiDAR (or RGB-D camera) on the right side to obtain direct
observations. This is the most reliable solution, as it eliminates
extrapolation entirely and restores the original ``rich viewpoints''
assumption.
\end{answerbox}

\end{document}
